\documentclass[10pt]{article}

% ---------- Document Identity (update these only) ----------
\newcommand{\DocStage}{}
\newcommand{\DocTitle}{Your Road to Tier One \textbf{SOF}}
\newcommand{\ProgramName}{182-Week Civilian-to-Selection Readiness Pipeline}
\newcommand{\ProgramSubtitle}{}

% ---------- Page ----------
\usepackage[legalpaper,left=0.5in,right=0.5in,top=0.6in,bottom=0.45in]{geometry}

% ---------- Typography ----------
\usepackage[T1]{fontenc}
\usepackage[utf8]{inputenc}
\usepackage{libertinus}
\usepackage{microtype}
\usepackage{parskip}
\usepackage{ragged2e}
\usepackage{textcomp}
\AtBeginDocument{\color{black}}
\AtBeginDocument{\pagecolor{white}}
\AtBeginDocument{\justifying}

% ---------- Landscape pages ----------
\usepackage{pdflscape}

% ---------- Color ----------
\usepackage[table]{xcolor}
\definecolor{StageBlue}{HTML}{1F4E79}
\definecolor{StageBlueDark}{HTML}{163A5A}
\definecolor{LightGray}{HTML}{F2F4F7}
\definecolor{MidGray}{HTML}{D9DEE5}
\definecolor{RuleGray}{HTML}{C7CED8}
\definecolor{HeaderInk}{HTML}{000000}
\definecolor{Ink}{HTML}{000000}

% ---------- Utility ----------
\usepackage{enumitem}
\sloppy
\setlength{\emergencystretch}{2em}

% ---------- Boxes / UI ----------
\usepackage[most]{tcolorbox}

\tcbset{
  charterTitle/.style={
    enhanced,
    colback=StageBlue!4,
    colframe=StageBlue,
    boxrule=1.25pt,
    arc=3pt,
    left=16pt,right=16pt,top=14pt,bottom=14pt,
    borderline north={3.2pt}{0pt}{StageBlue},
    borderline west={4pt}{0pt}{StageBlue},
    borderline south={1.25pt}{0pt}{StageBlue},
    drop shadow={black!25!white},
  },
  charterRule/.style={
    enhanced,
    colback=LightGray,
    colframe=RuleGray,
    boxrule=0.85pt,
    arc=3pt,
    left=12pt,right=12pt,top=10pt,bottom=10pt,
    borderline west={3pt}{0pt}{StageBlue},
  },
  charterBadge/.style={
    enhanced,
    colback=StageBlue,
    coltext=white,
    colframe=StageBlueDark,
    boxrule=0.6pt,
    arc=2pt,
    left=6pt,right=6pt,top=3pt,bottom=3pt,
  }
}
\newtcbox{\Badge}{nobeforeafter,charterBadge}

% ---------- Lists ----------
\setlist[itemize]{leftmargin=1.5em, itemsep=0.25em, topsep=0.25em}
\setlist[enumerate]{leftmargin=1.8em, itemsep=0.25em, topsep=0.25em}

% ---------- TOC ----------
\usepackage{tocloft}
\setcounter{tocdepth}{3}
\setlength{\cftbeforesecskip}{3pt}
\setlength{\cftbeforesubsecskip}{2pt}
\setlength{\cftbeforesubsubsecskip}{0pt}
\renewcommand{\cftsecfont}{\bfseries}
\renewcommand{\cftsecpagefont}{\bfseries}
\renewcommand{\cftsubsecfont}{\normalfont}
\renewcommand{\cftsubsecpagefont}{\normalfont}
\renewcommand{\cftsubsubsecfont}{\normalfont}
\renewcommand{\cftsubsubsecpagefont}{\normalfont}
\setlength{\cftsecindent}{0pt}
\setlength{\cftsubsecindent}{1.25em}
\setlength{\cftsubsubsecindent}{2.8em}
\setlength{\cftsecnumwidth}{2.1em}
\setlength{\cftsubsecnumwidth}{2.8em}
\setlength{\cftsubsubsecnumwidth}{3.6em}

% Remove the built-in TOC heading so only the custom heading is shown
\makeatletter
\renewcommand{\tableofcontents}{\@starttoc{toc}}
\makeatother

% ---------- Header/Footer: page number only ----------
\usepackage{fancyhdr}
\pagestyle{fancy}
\fancyhf{}
\renewcommand{\headrulewidth}{0pt}
\renewcommand{\footrulewidth}{0pt}
\fancyfoot[C]{\thepage}

% ---------- Section formatting ----------
\usepackage{titlesec}

\titleformat{\section}
  {\Large\bfseries\color{Ink}}
  {\thesection}{0.6em}{}
  [\vspace{0.25em}\textcolor{StageBlue}{\rule{\linewidth}{0.8pt}}]

\titleformat{\subsection}
  {\large\bfseries\color{Ink}}
  {\thesubsection}{0.6em}{}

\titleformat{\subsubsection}
  {\normalsize\bfseries\color{Ink}}
  {\thesubsubsection}{0.6em}{}

\titlespacing*{\section}{0pt}{0.95em}{0.35em}
\titlespacing*{\subsection}{0pt}{0.65em}{0.25em}
\titlespacing*{\subsubsection}{0pt}{0.55em}{0.2em}

% ---------- Table engine ----------
\usepackage{tabularray}
\UseTblrLibrary{booktabs}

% ---------- tabularray: GLOBAL table treatment ----------
\SetTblrInner{
  cells = {fg=black, bg=white, valign=t},
}

% ---------- Header helpers ----------
\newcommand{\Hdr}[1]{\shortstack[c]{\strut #1\strut}}

% ---------- Lines ----------
\newcommand{\TitleLine}{\textcolor{StageBlue}{\rule{0.98\linewidth}{0.95pt}}}
\newcommand{\SubTitleLine}{\textcolor{RuleGray}{\rule{0.98\linewidth}{0.65pt}}}

% ---------- Continued on next page ----------
\newcommand{\ContinuedOnNextPageInline}{%
  \par
  \vfill
  \noindent\textcolor{RuleGray}{\rule{\linewidth}{1pt}}\vspace{-2pt}
  \noindent\hfill\textit{Continued on next page}\par\vspace{-10pt}
  \noindent\textcolor{RuleGray}{\rule{\linewidth}{1pt}}
  \clearpage
}

% ---------- PDF hyperlinks ----------
\usepackage[hidelinks]{hyperref}
\hypersetup{
  colorlinks=false,
  pdfborder={0 0 0},
  linktoc=all,
  pdftitle={\DocTitle},
  pdfauthor={\ProgramName}
}

% =========================================================
\begin{document}

\section*{Stage 12.D — Emergency Stop Rules and Escalation Flow}
\phantomsection
\addcontentsline{toc}{section}{Stage 12.D — Emergency Stop Rules and Escalation Flow}

\subsection*{12.D.1 Authority and Boundary Statement}
\phantomsection
\addcontentsline{toc}{subsection}{12.D.1 Authority and Boundary Statement}

\begin{itemize}
\item This appendix aligns to Stage 0.5 safety and legal boundaries and Stage 1.F stop-rule doctrine and \textbf{MUST} be applied as written.
\item This appendix is non-medical governance only:
  \begin{itemize}
  \item no diagnosis,
  \item no treatment,
  \item escalation and documentation posture only.
  \end{itemize}
\item Safety overrides schedule, performance goals, and evidence collection convenience.
\item \textbf{STOP} \textbf{NOW} is mandatory when a stop trigger is present.
\item Evidence captured during a stop-rule breach is inadmissible for gate use and \textbf{MUST} NOT support \textbf{ADVANCE}.
\item Any impacted test window \textbf{MUST} follow Stage 3 reslot posture and no-stacking discipline.
\end{itemize}

\subsection*{12.D.2 Immediate Stop Triggers — Quick Reference}
\phantomsection
\addcontentsline{toc}{subsection}{12.D.2 Immediate Stop Triggers — Quick Reference}

\begingroup
\scriptsize
\setlength{\tabcolsep}{2pt}
\renewcommand{\arraystretch}{1.12}

\begin{longtblr}{
  width=\linewidth,
  rowhead=1,
  colspec = {
    X[l,1.55]
    X[l,2.35]
    X[l,2.10]
    X[l,4.00]
  },
  hlines,
  vlines,
  cells = {hsep=2pt, vsep=6pt, halign=l, valign=t},
  row{odd} = {bg=white},
  row{even} = {bg=LightGray},
  row{1} = {bg=StageBlue!15, fg=black, font=\bfseries, halign=l, valign=m},
}
\Hdr{Trigger Category} &
\Hdr{Examples} &
\Hdr{Immediate Action} &
\Hdr{Minimum Documentation} \\
Cardiovascular red flags &
Chest pain/pressure; fainting/syncope; near-syncope; disproportionate shortness of breath; palpitations with distress; collapse &
\textbf{STOP} \textbf{NOW}; cease activity; sit/lie down safely; EMS first if immediate danger; otherwise contact Medical &
Timestamp ISO 8601 with timezone offset; Phase and Week and Session \# or Test Event marker; Environment Unit \textbf{ID} (\textbf{ROUTE}-\#\#/\textbf{POOL}-\#\#/\textbf{MACH}-\#\#); Activity being performed; Trigger category; Observed signs and symptoms non-diagnostic; Immediate action taken; Session validity tag \textbf{VALID}/\textbf{MODIFIED}/\textbf{INVALID} plus reason code posture OR Test validity tag \textbf{VALID}/\textbf{INVALID} plus reason code posture; EMS or Medical contacted Y/N; Who was notified and when \\
Neurologic red flags &
New neurologic symptoms; severe headache with concerning features; confusion/disorientation; loss of coordination; vision changes; seizure-like activity &
\textbf{STOP} \textbf{NOW}; cease activity; stabilize; EMS first if immediate danger; otherwise contact Medical &
Timestamp ISO 8601 with timezone offset; Phase/Week + Session \# or Test marker; Environment Unit \textbf{ID}; Activity; Trigger category; Observed signs non-diagnostic; Immediate action; Validity tag + reason code; EMS/Medical contacted Y/N; Notifications \\
Heat illness signs &
Dizziness; nausea/vomiting; chills/goosebumps in heat; severe headache; confusion; cessation of sweating with distress; inability to continue safely &
\textbf{STOP} \textbf{NOW}; move to shade/indoors; cool down posture; EMS first if severe/confusion/collapse; otherwise contact Medical &
Timestamp ISO 8601 with timezone offset; Phase/Week + Session \# or Test marker; Environment Unit \textbf{ID}; Activity; Trigger category; Observed signs; Immediate action; Validity tag + reason code; EMS/Medical contacted Y/N; Notifications \\
Acute injury with instability or swelling &
Joint instability; rapid swelling; audible pop with immediate functional loss; inability to bear weight safely &
\textbf{STOP} \textbf{NOW}; cease load-bearing; stabilize; contact Medical &
Timestamp ISO 8601 with timezone offset; Phase/Week + Session \# or Test marker; Environment Unit \textbf{ID}; Activity; Trigger category; Observed signs; Immediate action; Validity tag + reason code; Medical contacted Y/N; Notifications \\
Gait-altering pain or mechanics change &
Limp; compensation; altered stride; unstable joint position; pain that changes mechanics &
\textbf{STOP} \textbf{NOW} for the exposure; substitute to lower-risk modality if safe; contact Medical if concerning &
Timestamp ISO 8601 with timezone offset; Phase/Week + Session \# or Test marker; Environment Unit \textbf{ID}; Activity; Trigger category; Observed signs; Immediate action; Validity tag + reason code; Medical contacted Y/N; Notifications \\
Water safety violations &
Any unsupervised underwater/breath-hold attempt; distress in water; disorientation/panic; cramp; near-drowning indicators &
\textbf{STOP} \textbf{NOW}; exit water immediately; lifeguard engaged if present; account for buddy; no re-entry same day; EMS if distress/near-drowning indicators; contact Medical &
Timestamp ISO 8601 with timezone offset; Phase/Week + Session \# or Test marker; Environment Unit \textbf{ID} (\textbf{POOL}-\#\#); Activity; Trigger category; Observed signs; Immediate action; Validity tag + reason code; EMS/Medical contacted Y/N; Notifications \\
Mental and behavioral stop triggers &
Acute panic/dissociation; suicidal ideation/self-harm statements; inability to safely self-regulate in-session; psychotic or manic symptoms &
\textbf{STOP} \textbf{NOW}; remove from hazard; do not continue; escalate to professional support posture; contact Medical when indicated &
Timestamp ISO 8601 with timezone offset; Phase/Week + Session \# or Test marker; Environment Unit \textbf{ID}; Activity; Trigger category; Observed signs non-diagnostic; Immediate action; Validity tag + reason code; EMS/Medical contacted Y/N; Notifications \\
\end{longtblr}

\endgroup

\newpage

\subsection*{12.D.3 Escalation Flow — Decision Tree}
\phantomsection
\addcontentsline{toc}{subsection}{12.D.3 Escalation Flow — Decision Tree}

\begin{enumerate}
\item Stop
  \begin{itemize}
  \item \textbf{STOP} \textbf{NOW} immediately.
  \item Terminate the exposure (session element or test attempt) without negotiation.
  \item If a test attempt is in progress, treat the test as \textbf{INVALID} unless all protocol completion criteria were already satisfied prior to the stop trigger.
  \end{itemize}

\item Stabilize and Exit Environment
  \begin{itemize}
  \item Move to a safer setting:
    \begin{itemize}
    \item stop movement; sit/lie down as appropriate; avoid hazards.
    \end{itemize}
  \item If water-related:
    \begin{itemize}
    \item exit water immediately,
    \item lifeguard engaged if present,
    \item buddy accounted for,
    \item no re-entry the same day,
    \item log as \textbf{SAFETY},
    \item default to \textbf{MEDICAL} \textbf{HOLD} if distress or near-drowning indicators occurred.
    \end{itemize}
  \end{itemize}

\item Notify and Escalate by Role
  \begin{itemize}
  \item Emergency-services-first pathway (life-threatening or severe red flags)
    \begin{itemize}
    \item EMS first when immediate danger is present.
    \item Then notify Medical.
    \item Then notify Coach Lead.
    \item Then notify others as needed (Equipment Lead, Trainee support contact) consistent with local policy.
    \end{itemize}
  \item Non-life-threatening but concerning pathway
    \begin{itemize}
    \item Medical first (non-diagnostic evaluation posture).
    \item Then notify Coach Lead.
    \item Then notify others as needed.
    \end{itemize}
  \item Mental and behavioral hard response pathway
    \begin{itemize}
    \item \textbf{STOP} \textbf{NOW}.
    \item Remove from hazard.
    \item Do not continue the session or test.
    \item Escalate to professional support posture.
    \item Document as \textbf{SAFETY}.
    \item Apply \textbf{MEDICAL} \textbf{HOLD} posture when safety is compromised.
    \end{itemize}
  \end{itemize}

\item Document Validity and Context
  \begin{itemize}
  \item Record immediately (same-day) using Stage 2.E posture:
    \begin{itemize}
    \item Timestamp ISO 8601 with timezone offset.
    \item Phase, Phase Week, Session \# (1–6) or Test Event marker.
    \item Environment Unit \textbf{ID}: \textbf{ROUTE}-\#\# or \textbf{POOL}-\#\# or \textbf{MACH}-\#\# as applicable.
    \item Activity being performed and the exposure stage (warm-up/main work/test attempt).
    \item Observed signs and symptoms (non-diagnostic) and trigger category.
    \item Immediate actions taken and escalation contacts (EMS Y/N; Medical Y/N; who notified and when).
    \item Validity tag and reason code:
      \begin{itemize}
      \item session: \textbf{VALID} / \textbf{MODIFIED} / \textbf{INVALID} plus \textbf{SESSION-RC} family (\textbf{ENV}/PAIN/\textbf{EQUIP}/\textbf{TIME}/\textbf{ILL}/\textbf{SAFETY}),
      \item test: \textbf{VALID} / \textbf{INVALID} plus \textbf{TEST-RC} family (\textbf{ENV}/PAIN/\textbf{EQUIP}/\textbf{TIME}/\textbf{ILL}/\textbf{SAFETY}).
      \end{itemize}
    \item Tool context when applicable (Timer/Scale/Tape identity; machine settings if relevant).
    \item Reslot posture note if a protected test window is affected:
      \begin{itemize}
      \item “Reslot to next deload/test week per Stage 3; no stacking make-ups.”
      \end{itemize}
    \end{itemize}
  \end{itemize}

\item Return-to-Training Gate High-Level
  \begin{itemize}
  \item Return-to-training is governed by safety and admissibility:
    \begin{itemize}
    \item do not resume maximal efforts or gate tests until safe and cleared under appropriate posture.
    \end{itemize}
  \item Gate decisions:
    \begin{itemize}
    \item any missing admissible evidence defaults to \textbf{HOLD} or \textbf{MEDICAL} \textbf{HOLD} as appropriate.
    \item reslot affected tests to the next protected window; do not stack make-ups.
    \end{itemize}
  \end{itemize}
\end{enumerate}

\newpage

\subsection*{12.D.4 Post-Event Re-Entry Rules — High-Level}
\phantomsection
\addcontentsline{toc}{subsection}{12.D.4 Post-Event Re-Entry Rules — High-Level}

\begin{itemize}
\item Cancel or postpone testing when stop triggers occur.
\item Reslot impacted tests to the next deload/test week per Stage 3.
\item Do not stack make-up tests.
\item Default gate outcome when admissible evidence is missing:
  \begin{itemize}
  \item \textbf{HOLD} or \textbf{MEDICAL} \textbf{HOLD} as appropriate.
  \end{itemize}
\end{itemize}

Default postures by category (non-medical):

\begin{itemize}
\item Cardiovascular red flags, neurologic red flags, suspected heat illness, loss of consciousness, water distress events:
  \begin{itemize}
  \item default \textbf{MEDICAL} \textbf{HOLD},
  \item clinician clearance required before resuming maximal efforts or gate tests.
  \end{itemize}
\item Gait-altering pain or mechanics change:
  \begin{itemize}
  \item default \textbf{HOLD},
  \item regress exposure and substitute lower-risk modalities,
  \item reslot impacted tests.
  \end{itemize}
\item Equipment or environment safety event:
  \begin{itemize}
  \item default \textbf{HOLD} until controls verified and environment/tooling validity restored (correct Environment Unit \textbf{ID} and Tool Standard posture available).
  \end{itemize}
\item Mental and behavioral stop trigger events:
  \begin{itemize}
  \item \textbf{MEDICAL} \textbf{HOLD} posture when safety is compromised,
  \item professional support required,
  \item no gate testing until stabilized.
  \end{itemize}
\end{itemize}

\end{document}